\scanpage{8}
\mspage{7}

This structure does not change if
\deleted{\unsure{of}}one replaces the order by the opposite order,
i.e. it depends only on the
biorder.

\deleted{Tame subset} \deleted{\unsure{Segments|Segment}}

% slanted
\margin{redo, without restricting oneself to a line : closed sub-segment piece of a segment piece. It is a segment piece \unsure{for|for} the induced biorder. On the other hand, the boundary of a segment piece is defined As some}

\unsure{Segment pieces|Segment piece} of the line : subsets
of the form

\[
\left\{\begin{aligned}
 &[s,t] &&s\leq t\\
 &\mathopen{]}-\infty,s],\quad\text{or}\quad [s,+\infty\mathclose{[}\\
 &\mathcal L=\mathopen{]}-\infty,+\infty\mathclose{[}
\end{aligned}\right.
\]

Boundary of the segment piece : $\{s,t\}$ in the
first case, $\{s\}$ in the second, $\varnothing$ \inserted{third}

Two segment pieces are said to be in
tame relative position, if their
boundaries : are so.

\uline{Tame} \uline{subset} :
\inserted{finite} union of segment pieces \inserted{in a position} mutually
tame, and disjoint. \deleted{Thus} By
definition, its boundary is the union
(disjoint) of the boundaries.

% paper watermark; clipped at the two sheet edges
% the paper's watermark (EXTRA STRONG), not his text
